\documentclass[10pt,a4paper]{amsart}
\usepackage{amsmath,amssymb,amsthm,graphicx}
\usepackage{hyperref}
\usepackage{tikz}
%\usepackage{xcolor}
\usepackage{pict2e}
\let\arc\relax
%\usepackage{harpoon}
\usepackage{curves}
%\usepackage{chngcntr}
%\usepackage{apptools}
\usepackage{enumerate}

%\usepackage{verbatim}
%\usepackage[nospace,noadjust]{cite}
\usepackage{accents}
\usetikzlibrary{calc,decorations.markings}
\unitlength1cm

%%%%%%%%%%%%%%FONTS%%%%%%%%%%%%%%%%%%%%%%%%%%%%%%
\newcommand{\R}{\mathbb{R}}
\newcommand{\Z}{\mathbb{Z}}
\newcommand{\N}{\mathbb{N}}
\newcommand{\C}{\mathbb{C}}
\newcommand{\Id}{I}

%%%%%%%%%%%%%%MATHCAL%%%%%%%%%%%%%%%%%%%%%%%
\newcommand{\cK}{\mathcal{K}}
\newcommand{\cE}{\mathcal{E}}
\newcommand{\cR}{\mathcal{R}}
\newcommand{\cW}{\mathcal{W}}

%%%%%%%%%%%%%%%%%%ABBRS%%%%%%%%%%%%%%%%%%%%%%%%%%%%%

%%%%%%% OBJECTS SPECIFIC TO THIS PAPER%%%%%%%%%

\newcommand{\hw}{\widehat w}
\newcommand{\tw}{\widetilde w}

\newcommand{\bQ}{ Q}
\newcommand{\hQ}{{\widehat Q}}
\newcommand{\tQ}{{\widetilde Q}}

\newcommand{\sbQ}{Q^\star}
\newcommand{\shQ}{{\widehat Q^\star}}
\newcommand{\stQ}{{\widetilde Q^\star}}

\newcommand{\smu}{\mu^\star}
\newcommand{\shmu}{\widehat \mu^\star}
\newcommand{\stmu}{\widetilde \mu^\star}

\newcommand{\bmu}{{\mu}}
\newcommand{\hmu}{{\widehat \mu}}
\newcommand{\tmu}{{\widetilde \mu}}


\newcommand{\bP}{P}
\newcommand{\hP}{{\widehat P}}
\newcommand{\tP}{{\widetilde P}}

\newcommand{\sbP}{P^\star}
\newcommand{\shP}{{\widehat P^\star}}
\newcommand{\stP}{{\widetilde P^\star}}

\newcommand{\hR}{{\widehat R}}

\newcommand{\bv}{v}
\newcommand{\hv}{{\widehat v}}
\newcommand{\tv}{{\widetilde v}}

\newcommand{\sbv}{v^\star}
\newcommand{\shv}{{\widehat v^\star}}
\newcommand{\stv}{{\widetilde v^\star}}

\newcommand{\stu}{{\widetilde u^\star}}

\newcommand{\hF}{\widehat F}
\newcommand{\hW}{\widehat W}

\newcommand{\Fw}{\frac{F^2}{w}}
\newcommand{\hFw}{\frac{\hF^2}{\hw}}

\newcommand{\Cau}{\mathcal C_{\Sigma}}
\newcommand{\Cauv}{\mathcal C_{v}}
\newcommand{\Cauhv}{\mathcal C_{\hv}}
\newcommand{\Cautv}{\mathcal C_{\tv}}

\newcommand{\sCau}{\mathcal C_{\Sigma^\star}}
\newcommand{\sCauv}{\mathcal C_{\bv^\star}}
\newcommand{\sCauhv}{\mathcal C_{\hv^\star}}
\newcommand{\sCautv}{\mathcal C_{\stv}}

\newcommand{\Ltwo}{L^2(\Sigma)}
\newcommand{\sLtwo}{L^2(\Sigma^\star)}

\newcommand{\Sigl}{\Sigma^\ell}
\newcommand{\Sigr}{\Sigma^r}

%%%%%%%%%%%%%%%%OTHER SHORTCUTS%%%%%%%%%%%
\newcommand{\nn}{\nonumber}
\newcommand{\ul}{\underline}
\newcommand{\ol}{\overline}
\newcommand{\pa}{\partial}
\newcommand{\ti}{\tilde}
\newcommand{\wti}{\widetilde}
\newcommand{\wha}{\widehat}
\newcommand{\norm}[1]{\lVert#1 \rVert}
\newcommand{\abs}[1]{\lvert#1 \rvert}
%\newcommand{\Id}{\mathbb{I}}
\newcommand{\I}{\mathrm{i}}
\newcommand{\E}{\mathrm{e}}
\newcommand{\re}{\mathop{\mathrm{Re}}}
\newcommand{\im}{\mathop{\mathrm{Im}}}
\newcommand{\eps}{\varepsilon}
\newcommand{\dist}{\mathop{\mathrm{dist}}}

\newtheorem{Theorem}{Theorem}[section]
\newtheorem{Corollary}[Theorem]{Corollary}
\newtheorem{Lemma}[Theorem]{Lemma}
\newtheorem{Proposition}[Theorem]{Proposition}
 { \theoremstyle{definition}
\newtheorem{Definition}[Theorem]{Definition}
\newtheorem{Note}[Theorem]{Note}
\newtheorem{Example}[Theorem]{Example}
\newtheorem{Remark}[Theorem]{Remark} }

\numberwithin{equation}{section}
\usepackage{cleveref}
\crefname{originalequation}{original equation}{original equations}
\Crefname{originalequation}{Original equation}{Original equations}
\makeatletter
\newlabel{BesselAsym0}{{internal}{source}{}{sourcepointer-1}{}}
\newlabel{BesselAsym0@cref}{{[originalequation][0][]internal}{[1][0][]source}{}{}{}}
\makeatother
\crefname{originalequation}{original equation}{original equations}
\Crefname{originalequation}{Original equation}{Original equations}
\makeatletter
\newlabel{BesselAsym1}{{internal}{source}{}{sourcepointer-2}{}}
\newlabel{BesselAsym1@cref}{{[originalequation][0][]internal}{[1][0][]source}{}{}{}}
\makeatother
\crefname{originalequation}{original equation}{original equations}
\Crefname{originalequation}{Original equation}{Original equations}
\makeatletter
\newlabel{BesselInfinity}{{internal}{source}{}{sourcepointer-3}{}}
\newlabel{BesselInfinity@cref}{{[originalequation][0][]internal}{[1][0][]source}{}{}{}}
\makeatother
\crefname{originalequation}{original equation}{original equations}
\Crefname{originalequation}{Original equation}{Original equations}
\makeatletter
\newlabel{defXi}{{internal}{source}{}{sourcepointer-4}{}}
\newlabel{defXi@cref}{{[originalequation][0][]internal}{[1][0][]source}{}{}{}}
\makeatother
\title{Jump, matching and growth properties of the three local parametrices}
\author{Deift, Percy, Piorkowski, Mateusz}
\date{doi:10.3842/SIGMA.2024.004; CC BY 4.0}
\begin{document}
\maketitle
\setcounter{section}{5}
\setcounter{Theorem}{1}
\noindent\textit{Editorial note.} All delivered mathematical text is the author's own, copied from the pinned SIGMA article source; no statement or proof step was rewritten, reordered or removed. The counted statement is the labeled lemma environment at source lines 1363-1390 and the counted proof is the proof environment at source lines 1391-1457.\par\medskip
\noindent\textit{Editorial note.} The proof attributes the construction of the parametrices and the first two properties of the statement to a published section of the cited literature, which is declared as a standard prerequisite of this card; the remaining properties are derived in the delivered text.\par\medskip
\noindent\textit{Editorial note.} Two displays of earlier propositions are delivered as context spans because the counted proof refers to them by labels that the external-reference declaration of the delivery format cannot carry; the context spans themselves contain no reference to labels outside the delivered text.\par\medskip
\noindent\textit{Editorial note.} The proof cites published references of the author's bibliography by key; the cited entries are copied line for line from the author's own bibliography into a retained bibliography fragment, and hyperref is added because the publisher class, not the author preamble, supplied the link command.\par\medskip
\noindent\textit{Editorial note.} The retained preamble of this unit was generated mechanically for this delivery: it copies the author's own package, numbering, macro and theorem declaration block from the pinned source (source lines 9-133), with the publisher class and page setup lines removed.\par\medskip
\noindent\textit{Allowed standard prerequisites (editorial).}\begin{itemize}
\item The construction of the local parametrices and their basic properties. published section cited by the author at proof line 1392. Supplies the jump relations on the contour and the matching condition on the boundary, which the author attributes to that published section instead of rederiving them.
\item The Bessel asymptotics used to expand the model parametrices. same article, the displayed asymptotics at source lines 1288, 1294 and 1305. Converts the model problem solved by Bessel functions into the two sided growth bound of the statement near the endpoint minus one.
\item The definition of the phase function used in the local analysis. same article, the display at source line 1192. Enters the explicit form of the model parametrix whose behaviour is estimated.
\item The asymptotics of the squared Szegoe quotients at the endpoint minus one. same article, the displays delivered with this card as context spans. Provides the two exponential representations that are substituted into the model parametrix before the growth bound is taken.
\end{itemize}\medskip
\ifdefined\hypertarget\hypertarget{sourcepointer-1}{}\fi
\noindent\textit{Source reference.} Original equation internal: sources/sigma-2024-004-recurrence-coefficients-for-orthogonal-polyn/pinned-original-source.tex, source lines 1288--1292. The counted statement or proof refers to that the author's displayed Bessel asymptotics used in the first regime; it is the author's own earlier display, definition or statement in the same article and the whole argument of the counted span is retained verbatim.\par\medskip
\ifdefined\hypertarget\hypertarget{sourcepointer-2}{}\fi
\noindent\textit{Source reference.} Original equation internal: sources/sigma-2024-004-recurrence-coefficients-for-orthogonal-polyn/pinned-original-source.tex, source lines 1294--1298. The counted statement or proof refers to that the author's displayed Bessel asymptotics used in the second regime; it is the author's own earlier display, definition or statement in the same article and the whole argument of the counted span is retained verbatim.\par\medskip
\ifdefined\hypertarget\hypertarget{sourcepointer-3}{}\fi
\noindent\textit{Source reference.} Original equation internal: sources/sigma-2024-004-recurrence-coefficients-for-orthogonal-polyn/pinned-original-source.tex, source lines 1305--1309. The counted statement or proof refers to that the author's displayed Bessel asymptotics at infinity; it is the author's own earlier display, definition or statement in the same article and the whole argument of the counted span is retained verbatim.\par\medskip
\ifdefined\hypertarget\hypertarget{sourcepointer-4}{}\fi
\noindent\textit{Source reference.} Original equation internal: sources/sigma-2024-004-recurrence-coefficients-for-orthogonal-polyn/pinned-original-source.tex, source lines 1192--1196. The counted statement or proof refers to that the author's displayed definition of the phase function; it is the author's own earlier display, definition or statement in the same article and the whole argument of the counted span is retained verbatim.\par\medskip
\begin{Lemma}\label{PropP}
The matrix-valued functions $\bP$, $\hP$ and $\tP$ defined in $U \setminus \Sigma$, satisfy the following conditions:
\begin{enumerate}[$(i)$]\itemsep=0pt
 \item For $s \in U\cap \Sigma$,
 \begin{align} \label{ParametrixJump}
 \bP_+(s) = \bP_-(s) \bv(s), \qquad
 \hP_+(s) = \hP_-(s) \hv(s), \qquad
 \tP_+(s) = \tP_-(s) \tv(s).
 \end{align}
 \item For $s \in \partial U$ $($see \emph{\cite[equation~(6.41)]{KMVV}}$)$,
 \begin{align}\label{ParaMatchingCond}
 \bP(s), \hP(s), \tP(s) = N(s) + O\big(n^{-1}\big).
 \end{align}
\item For $z \in U$, we have uniformly $($see \emph{\cite[equation~(7.10)]{KMVV}}$)$
\begin{align} \label{ParaConv}
 \tQ(z) \big[\tP(z)\big]^{-1} = \Id + O\big(n^{-1}\big), \qquad \hQ(z) \big[\hP(z)\big]^{-1} = \Id + O\big(n^{-1}\big).
\end{align}
\item For $z \in U$,
\begin{align}\label{ParaBehaviour}
 \bP(z), \hP(z) = O\big(\max\big\lbrace |z+1|^{-1/4}, n^{-1/2} |z+1|^{-1/2}\big\rbrace\big),
\end{align}
and
\begin{align}\label{LegendreBehaviour}
 \tP(z) = O\big(|z+1|^{-1/4}\big)
\end{align}
uniformly as $n \to \infty$.
\end{enumerate}
\end{Lemma}

\begin{proof}
A detailed derivation of the local parametrices can be found in \cite[Section~6]{KMVV} together with a proof of properties (i), (ii), for property (iii) see \cite[Section~7]{KMVV}. Note that while the weight function $w$ does not fall into the class of weight functions considered in \cite{KMVV} due to the logarithmic singularity at $z = +1$, the local construction and estimation of the left parametrix near $z = -1$ found therein remains unchanged.

Regarding point (iv), we start with the properties of $\bP$ and $\hP$. Noting that $E(z)$ and~$\widehat E(z)$ are holomorphic, $n$-independent and have unit determinants, hence it is enough to consider $E^{-1}(z)\bP(z)$ and $\widehat E^{-1}(z)\hP(z)$ instead of $\bP(z)$ and $\hP(z)$ in \eqref{ParaBehaviour}. Analogously, it follows from~\eqref{F^2/w at -1} and \eqref{hat F^2/w at -1} that \smash{$\big(\frac{F}{W}\big)^{\sigma_3}$} and \smash{$\big(\frac{\hF}{\widehat W}\big)^{\sigma_3}$} are $n$-independent and bounded in a neighbourhood of $z = -1$, hence they also do not contribute in \eqref{ParaBehaviour}.

It remains to study
\begin{align*}
 (2\pi n)^{\sigma_3/2}\Psi_1\big(n^2 \xi(z)\big)[-\phi(z)]^{-n\sigma_3}
\end{align*}
which is equal to both $E^{-1}(z)\bP(z)\big(\frac{F}{W}\big)^{-\sigma_3}$ and $\widehat E^{-1}(z)\hP(z)\big(\frac{\hF}{\widehat W}\big)^{-\sigma_3}$. It follows from the definition of $\xi$ in \eqref{defXi} that \smash{$[-\phi(z)]^{-n\sigma_3} =\E^{-2\sqrt{n^2 \xi(z)}\sigma_3}$} where the square root has a branch cut along~$z > -1$. Writing $\zeta=n^2 \xi(z)$ and assuming $|z+1| \gtrsim O\big(n^{-2}\big)$, we see that $|\zeta| \gtrsim O(1)$ and using the estimates in \eqref{BesselInfinity} we conclude that
\begin{align}\nonumber
& (2\pi n)^{\sigma_3/2}\Psi_1\big(n^2 \xi(z)\big)[-\phi(z)]^{-n\sigma_3} \\
&\quad= (2\pi n)^{\sigma_3/2}\begin{pmatrix}
 O\big(|\zeta|^{-1/4}\big) & O\big(|\zeta|^{-1/4}\big)
 \\
 O\big(|\zeta|^{1/4}\big) & O\big(|\zeta|^{1/4}\big)
 \end{pmatrix}\nonumber
 \\
 &\quad=
 \begin{pmatrix}
 O\big(|z+1|^{-1/4}\big) & O\big(|z+1|^{-1/4}\big)
 \\
 O\big(|z+1|^{1/4}\big) & O\big(|z+1|^{1/4}\big)
 \end{pmatrix}, \qquad |z+1| \gtrsim O\big(n^{-2}\big).\label{geqn2}
\end{align}
For $|z+1| \lesssim O\big(n^{-2}\big)$, we use the estimates \eqref{BesselAsym1} instead to conclude
\begin{align}
 &(2\pi n)^{\sigma_3/2}\Psi_1\big(n^2 \xi(z)\big)[-\phi(z)]^{-n\sigma_3} \nonumber\\
 &\quad= (2\pi n)^{\sigma_3/2}\begin{pmatrix}
 O\big(|\zeta|^{-1/2}\big) & O\big(|\zeta|^{-1/2}\big)
 \\
 O\big(|\zeta|^{-1/2}\big) & O\big(|\zeta|^{-1/2}\big)
 \end{pmatrix}\nonumber\\
&\quad= \begin{pmatrix}
 O\big(n^{-1/2}|z+1|^{-1/2}\big) & O\big(n^{-1/2}|z+1|^{-1/2}\big)
 \\
 O\big(n^{-3/2}|z+1|^{-1/2}\big) & O\big(n^{-3/2}|z+1|^{-1/2}\big)
 \end{pmatrix}, \qquad z \in |z+1| \lesssim O\big(n^{-2}\big).\label{seqn2}
\end{align}
Note that in this case $[-\phi(z)]^{-n\sigma_3} = O(1)$, hence this term does not contribute. One checks that indeed for $|z+1| \sim n^{-2}$ the bounds in \eqref{geqn2} and \eqref{seqn2} are of the same order.

The proof of \eqref{LegendreBehaviour} works in a similar fashion. Again the holomorphic prefactor $\widetilde E$ can be ignored. For $|z+1|\gtrsim O\big(n^{-2}\big)$, we get as before
\begin{align*}
& (2\pi n)^{\sigma_3/2}\Psi_0\big(n^2 \xi(z)\big)[-\phi(z)]^{-n\sigma_3}
 \\
 &\quad=\begin{pmatrix}
 O\big(|z+1|^{-1/4}\big) & O\big(|z+1|^{-1/4}\big)
 \\
 O\big((|z+1|^{1/4}\big) & O\big((|z+1|^{1/4}\big)
 \end{pmatrix}, \qquad |z+1| \gtrsim O\big(n^{-2}\big).
\end{align*}
However, for $|z+1|\lesssim O\big(n^{-2}\big)$, we get different asymptotics after applying \eqref{BesselAsym0}. We obtain%
\begin{align} \nonumber
 &(2\pi n)^{\sigma_3/2}\Psi_0\big(n^2 \xi(z)\big) [-\phi(z)]^{-n\sigma_3}
 \\\label{LegendreGrowth}
 &\quad= \begin{pmatrix}
 O\big(n^{1/2} \big|\log\big(n^2(z+1)\big)\big|\big) & O\big(n^{1/2} \big|\log \big(n^2(z+1)\big)\big|\big)
 \\
 O\big(n^{-1/2}\big) & O\big(n^{-1/2}\big)
 \end{pmatrix}, \quad |z+1|\lesssim O\big(n^{-2}\big).
\end{align}
Now observe that for $|z+1|\lesssim O\big(n^{-2}\big)$, we have trivially $n^2|z+1|\lesssim O(1)$ and thus
\begin{align*}
 \big|\log\big(n^2(z+1)\big)\big| \lesssim \big|n^2(z+1)\big|^{-1/4}.
\end{align*}
This estimate, together with $n^{-1/2} \lesssim |z+1|^{-1/4}$ \big(in fact $n^{1/2} \lesssim |z+1|^{-1/4}$ holds\big), implies that the matrix entries in \eqref{LegendreGrowth} can be bounded by $O\big(|z+1|^{-1/4}\big)$ uniformly as $n \to \infty$, showing~\eqref{LegendreBehaviour} and finishing the proof.
\end{proof}

\begin{align}\label{F^2/w at -1}
 \frac{F^2}{w}(z) = \phi(z)^{-1}\exp\bigl\lbrace-\big(z^2-1\big)^{1/2}\big(d_0+O(|z+1|)\big)\bigr\rbrace,
\end{align}
uniformly for any path $z\to-1$ in $\C \setminus [-1,\infty)$, where
\begin{align}\label{d0}
 d_0 = \frac{1}{2\pi\I} \int_\gamma \frac{\log\big(w(\zeta)/(1+\zeta)\big)}{(\zeta^2-1)^{1/2}}\frac{{\rm d}\zeta}{\zeta+1}
\end{align}

\begin{align}\label{hat F^2/w at -1}
 \frac{\hF^2}{\hw}(z) = \phi(z)^{-1}\exp\bigl\lbrace-\big(z^2-1\big)^{1/2}d_0\bigr\rbrace.
\end{align}

\begin{thebibliography}{99}
\bibitem{KMVV}
Kuijlaars A.B.J., McLaughlin K.T.-R., Van~Assche W., Vanlessen M., The
 {R}iemann--{H}ilbert approach to strong asymptotics for orthogonal
 polynomials on {$[-1,1]$}, \href{https://doi.org/10.1016/j.aim.2003.08.015}{\textit{Adv. Math.}} \textbf{188} (2004), 337--398,
 \href{https://arxiv.org/abs/math.CA/0111252}{arXiv:math.CA/0111252}.
\end{thebibliography}
\end{document}
